% AUTO-GENERATED by scripts/build_topics.py — do not edit.
% Edit topics/bodies/topic_27.tex or topics/preamble.tex instead.
%% Placeholder. Replace with the written chapter.
\textit{Източници:} PDF \texttt{27.pdf}, снимки \texttt{27ма тема}, \texttt{SA.pdf}, \texttt{SA-SI.pdf}

\subsection*{Анотация от конспекта}
Архитектури на софтуерни системи. a. Понятие за софтуерна архитектура. Структури и изгледи (structures and views) на архитектурата. Необходимост от софтуерни архитектури. Влияние на архитектурата върху проекта и организацията. b. Изисквания към качеството (нефункционални изисквания) на системата. c. Архитектурни стилове. Многослоен стил, Модел-изглед контролер (MVC), Поточни (Pipe-and-filter) архитектури. Архитектура ориентирана към услуги. d. Проектиране на софтуерната архитектура. Процес за проектиране. Избор на подходящи структури. Последователност на създаване на архитектурата. Тактики (архитектурни решения) за постигане на желаните качествени показатели. e. Документиране на софтуерната архитектура. Предназначение и основни елементи в документацията на архитектурата. Литература: [30], [52].
